\documentclass[a4paper]{scrartcl}
% \usepackage{tgheros}
% \usepackage[T1]{fontenc}
\usepackage{fontspec}
\setsansfont{TeX Gyre Heros}[Ligatures=TeX]
\usepackage[all]{nowidow}

\usepackage[sorting=none]{biblatex}
\addbibresource{references.bib}

\usepackage{graphicx}
\graphicspath{{./figure}}

\usepackage{caption}
\captionsetup{format=plain}

\usepackage{trimclip}
% \usepackage{xspace}
\usepackage{csquotes}
\usepackage{hyperref}
\usepackage{todonotes} % [disable]
\usepackage{amsmath}
\usepackage{siunitx}
% \newfontfamily{\lmr}{Latin Modern Roman}
% \sisetup{text-font-command=\lmr}
% \sisetup{output-prefix={\ensuremath{\mu}}}
\usepackage{amssymb}

\newcommand{\figref}[1]{Fig.\,\ref{#1}}


\begin{document}

% \begin{titlepage}
\begin{center}
\vspace*{1cm}

\LARGE
  \textbf{Digital Filters}

\vspace{.2cm}

\normalsize
  Physical Measuring Techniques FK7063

\vspace{.5cm}
{Michael Vogt}

June 2026
\end{center}

\vspace{1cm}



\section*{Introduction}
In this computational lab, we construct several digital filters using Matlab
and compare them to each other. We also analyze data from an EKG signal
and determine ways to remove noise from it.


% \end{titlepage}



\section{Tasks}
\subsection{Constructing a Second Order Stop-Band Filter}
We define the frequency response of a stop-band filter whose gain is $0$ at $f = f_0 = \SI{50}{\Hz}$
with a bandwidth of $b = \SI{4}{\Hz}$ at \SI{-3}{\dB} (corresponding to an attenuation of $1/\sqrt(2)$).
The sampling rate is $F_s = \SI{1000}{\Hz}$.

In the space of the z-transform, the frequency $f_0$ corresponds to $\exp(i\Omega_0)$ with
$\Omega_0 = 2\pi \cdot f_0/F_s$. To get the desired attenuation, we place a zero
at this point, i.\,e. we add a factor $(1- \exp(i\Omega_0))$ to the denominator.
Since real signals always contain negative frequencies in the same amount as corresponding positive frequencies,
we also need to attenuate \SI{-50}{\Hz} (corresponding to $-\Omega_0$) in the same way,
so we add a factor $(1 - \exp(-i\Omega 0))$.

Now we need to add the property that the filter's maximum gain is $1$ and that it has the desired bandwidth.
This is achieved by adding a pole (i.e. a factor in the denominator)
that is not on the unit circle but still close to the zero that we added.
It can be parametrized as $r \exp(i\Omega 0)$ (with $r$ real). 
For values of z that are far away, the distance to the pole and the zero will be roughly the same,
so modulus of the numerator and denominator are equal, leading to a gain of $1$.
When z approaches the pole and the zero, there comes a point when it is separated from the
zero by a distance $d \approx (1-r)$ and from the pole by a distance $\approx \sqrt{d+(1-r)} \approx \sqrt{2}(1-r)$.
the ratio of the numerator and denominator is then $1/\sqrt(2)$
% the bandwidth therefore corresponds to a distance of 2d around the unit circle, and d is related to r by d=1-r,
% where r is the distance of the pole from the origin. 



\section{Conclusion}




\vfill
\printbibliography

\end{document}